%============================================================================
% ARGUMENTATION DOCUMENT (POLICY LAYER) --- BODY ONLY
% Body-only version of argumentation-policy-EN.tex for inclusion by the suite
% assembler (policy-EN.tex). No preamble, no \begin{document}, no titlepage,
% no \tableofcontents and no \printbibliography --- those are provided by the
% parent dossier. The standalone argumentation-policy-EN.tex remains the
% authoritative source; this file mirrors its sectional Q&A body.
%============================================================================

%============================================================================
% READER'S GUIDE
%============================================================================
\section*{How to read this module}
\addcontentsline{toc}{section}{How to read this module}

This module anticipates the objections most frequently raised against the
policy-layer proposals of the dossier and answers each on the evidence. Three
box types structure every exchange:

\begin{description}[nosep,leftmargin=2em]
  \item[\textcolor{red!60!black}{\textbf{Question}}] a good-faith technical or
    governance objection, reformulated at its strongest before it is answered.
  \item[\textcolor{orange!70!black}{\textbf{Insidious framing}}] an objection
    that smuggles in a premise, a motive, or a rhetorical trap; the answer
    addresses the framing, not only the surface claim.
  \item[\textcolor{OliveGreen!60!black}{\textbf{Answer}}] an
    evidence-based response following the concession-then-rebuttal pattern:
    partial merit is conceded honestly, then the objection is answered across
    independent evidence layers.
\end{description}

\noindent\textbf{Evidence methodology.} Every substantive answer is built from
up to five independent layers, presented strongest-first and cross-corroborated:
\begin{enumerate}[nosep,leftmargin=2em]
  \item \keyword{Legal} --- binding instruments of the adopter's jurisdiction; in
    EU/CoE editions these are the EU Charter, \gls{ecthr} and \gls{cjeu} case law,
    \gls{gdpr}, \gls{nis2} and the \gls{aiact}, used here as the worked example;
  \item \keyword{Standards} --- normative frameworks (NIST CSF~2.0, NIST SP~800
    series, ISO/IEC~27001, \gls{cobit}~2019);
  \item \keyword{Sector} --- consensus of professional bodies and surveys
    (EDUCAUSE, \gls{geant}, Verizon DBIR, Sophos);
  \item \keyword{Empirical} --- quantified measurements and incident
    root-cause analysis;
  \item \keyword{Institutional} --- named universities and research
    organisations with publicly documented practice.
\end{enumerate}
\noindent A thesis is treated as robust only when at least three layers sustain
it independently; where a layer is contestable, the answer says so. The
institutional precedents named below (MIT, Stanford, Oxford, ETH~Z\"urich,
CERN, Michigan, and others) are cited as \emph{evidence}, not as models any
particular institution must replicate.

% ============================================================================
\section{Questions on network segmentation}
\label{sec:q-segmentation}
% ============================================================================

\epigraph{``For every complex problem there is an answer that is neat, simple, and wrong.''}{--- H.\,L.\ Mencken, \textit{Prejudices: Second Series}, 1920}

% --- Q1.1 ---
\begin{questionbox}[Q1.1 --- ``Segmentation increases complexity and therefore the attack surface'']
Network segmentation multiplies the number of zones, inter-zone boundaries, and firewall rules. More boundaries mean more configuration points, more potential misconfigurations, and therefore a \emph{larger} attack surface. A single, well-managed perimeter is simpler and therefore safer.
\end{questionbox}

\begin{answerbox}[Answer to Q1.1]
\textbf{Partial merit acknowledged:} segmentation does increase \emph{configuration complexity}. This is a genuine engineering concern, not a straw-man.

\medskip\noindent\textbf{However, the question conflates two distinct concepts:}
\begin{enumerate}[nosep]
    \item \emph{Configuration complexity} (the number of rules to manage) and
    \item \emph{Attack surface} (the set of points an attacker can exploit).
\end{enumerate}

Segmentation \emph{increases} the former but \emph{decreases} the latter. NIST SP~800-53 Rev.\,5~\cite{nist-sp800-53} explicitly recommends ``protecting internal boundaries and isolating components to reduce the attack surface and limit damage.'' The NIST Cybersecurity Framework~2.0~\cite{nist-csf} and ISO~27001:2022 control A.8.22~\cite{iso27001} both mandate network segregation precisely because it \emph{reduces} the blast radius of any single compromise.

\medskip\noindent\textbf{Empirical evidence:} documented university ransomware failures of 2024--2025 demonstrate that \emph{centralised, unsegmented perimeters amplify damage}. Where an institution runs a single large administrative perimeter, an attack can disrupt the entire institution for an extended period; conversely, a segmented architecture is designed to confine a breach to a single subnet, limiting the blast radius that a flat network exposes~\cite{rennes-cyber}. The Verizon 2024 DBIR~\cite{verizon-dbir2024} documents the prevalence of lateral movement in such breaches --- the propagation path that segmentation, per the NIST and ISO controls cited above, is designed to interrupt.

\medskip\noindent\textbf{Managing the configuration complexity:} the dossier's phased transition addresses this directly --- governance and charters first, a pilot on selected entities next, then gradual extension. \gls{sdn} tools and infrastructure-as-code automate rule management and reduce human error. The complexity argument is an argument for \emph{proper tooling and gradual rollout}, not for avoiding segmentation.
\end{answerbox}

\bigskip

% --- Q1.2 ---
\begin{questionbox}[Q1.2 --- ``A single perimeter is easier to monitor and defend'']
With all traffic passing through a single firewall, the security team has full visibility. Segmentation creates blind spots between zones.
\end{questionbox}

\begin{answerbox}[Answer to Q1.2]
This argument assumes that \emph{centralised inspection} equals \emph{comprehensive visibility}. The opposite is frequently true in practice:

\begin{itemize}[nosep]
    \item A single perimeter firewall handling research traffic at 10--100\,Gbps cannot perform deep packet inspection without creating bottlenecks --- a documented case measured a \textbf{five- to twelve-fold throughput degradation} when research traffic traversed an enterprise firewall~\cite{esnet-pennstate}.
    \item When performance degrades, researchers route around the firewall (Shadow IT), creating \emph{actual} blind spots far worse than managed inter-zone boundaries.
    \item In a segmented architecture, each zone has \emph{its own} monitoring (a Zeek \gls{ids} on the \gls{sciencedmz}, \gls{ids}/\gls{ips} on the administrative zone, per-zone logging baselines). Total observability is \emph{greater}, not lesser.
\end{itemize}

\noindent CERN explicitly operates zone-specific monitoring alongside backbone monitoring~\cite{cern-security-freedom}. The University of Michigan deploys per-zone security classifications and dedicated sensitive-data environments~\cite{umich-safecomputing}. These institutions have demonstrably \emph{better} visibility than a single-perimeter model, not worse.
\end{answerbox}

\bigskip

% --- Q1.3 ---
\begin{clarificationbox}[Q1.3 --- ``You just want less security to make researchers' lives easier'']
The real motivation behind this document is not improved security but reduced control. Researchers want to do whatever they want without oversight. This is a \emph{comfort} argument dressed up in security language.
\end{clarificationbox}

\begin{answerbox}[Answer to Q1.3]
The validity of a proposal is best assessed on its evidence and logic, independently of the motives attributed to its authors --- attacking the imputed motive rather than the argument is an \emph{ad hominem}. The proposals must be evaluated on their \emph{substance}.

\medskip\noindent\textbf{On the substance:} the dossier proposes \emph{more} security for sensitive data, not less. The administrative enclave is \emph{smaller, easier to harden and easier to monitor}. The security function gains:
\begin{itemize}[nosep]
    \item A reduced and well-defined perimeter for the most sensitive assets;
    \item Charters replacing informal exceptions with accountable, auditable autonomy;
    \item Zone-specific monitoring replacing a single, overwhelmed checkpoint;
    \item Evidence production for \gls{gdpr} Article~32 and \gls{nis2} compliance.
\end{itemize}

\noindent\textbf{The ``comfort'' framing is refuted by the proposal's own evidence:} the design is endorsed by NIST~\cite{nist-csf, nist-sp800-53}, ISO~27001~\cite{iso27001}, and EDUCAUSE~\cite{educause2024}, and is practiced by MIT, Stanford, Oxford, ETH~Z\"urich, CERN, and the 100-plus universities operating a \gls{sciencedmz}~\cite{esnet-history,nsf-cc}. If ``wanting less security'' were the true motive, these institutions and frameworks would not converge on the same model.

\medskip\noindent Constructive engagement with evidence-based proposals, as recommended by the \gls{cobit}~2019 process EDM05 (\emph{Ensured Stakeholder Engagement}), is more productive than debating motives~\cite{cobit}.
\end{answerbox}

\bigskip

% --- Q1.4 ---
\begin{questionbox}[Q1.4 --- ``The analogy with firebreaks is misleading --- a forest is not a network'']
The argument for segmentation leans on a forestry analogy (firebreaks). But networks are not forests: a misconfigured inter-zone rule \emph{is} a vulnerability, unlike a physical firebreak.
\end{questionbox}

\begin{answerbox}[Answer to Q1.4]
\textbf{Partial merit:} all analogies have limits, and pushing any metaphor too far weakens the argument. The firebreak analogy illustrates one principle only: \emph{compartmentalisation limits propagation}. It is not a complete model of network security.

\medskip\noindent However, the principle itself does not depend on the analogy. It is established by:
\begin{itemize}[nosep]
    \item NIST SP~800-207 (\gls{zt} Architecture)~\cite{nist-zt}: ``assume breach, limit blast radius'';
    \item NIST SP~800-53 Rev.\,5~\cite{nist-sp800-53}: ``isolate components to reduce attack surface'';
    \item The Verizon 2024 DBIR~\cite{verizon-dbir2024}: credential abuse and ransomware dominate the measured breach patterns --- the in-network progression that segmentation (per the NIST controls above) is designed to interrupt.
\end{itemize}

\noindent The risk of misconfigured rules is real but is an \emph{implementation risk}, not an argument against the principle. The phased transition (pilot first, then rollout) and use of \gls{sdn} automation directly mitigate this risk.

\smallskip\noindent\textit{\footnotesize Technical implementation: the segmentation contract, zone catalogue, and migration approach are developed in the standalone Technical Reference~\cite{ng-tech-ref-2026}.}
\end{answerbox}

\bigskip

% --- Q1.5 ---
\begin{questionbox}[Q1.5 --- ``You claim Zero Trust but exclude SASE/ZTNA and never map to NIST SP~800-207'']
You claim a Zero-Trust posture, yet you exclude SASE and ZTNA --- the two products the market equates with Zero Trust --- and never demonstrate Zero-Trust maturity against the actual standard, NIST SP~800-207.
\end{questionbox}

\begin{answerbox}[Answer to Q1.5]
\textbf{Excluding SASE/ZTNA from the structural foundation is a sovereignty choice, not a rejection of Zero Trust --- and the architecture maps directly to NIST SP~800-207~\cite{nist-zt}.} SASE and ZTNA are commercial \emph{packagings} of Zero-Trust ideas in which a supplier operates the data plane and terminates TLS; that is precisely the data- and technical-sovereignty loss the dossier argues against. Zero Trust is a set of principles, not a product, and the reference architecture realises them structurally:
\begin{enumerate}[nosep]
  \item \textbf{Resources, secured everywhere} (Tenets~1--2): every layer is a protected resource and communication is secured regardless of network location.
  \item \textbf{Identity-based, per-session access} (Tenets~3,~6): access is mediated by federated identity (\gls{sso}-capable identity providers, eduGAIN) and authorised per session, strictly enforced before access.
  \item \textbf{Dynamic policy} (Tenet~4): the policy plane decides on identity, role, posture and environment; zones default-deny.
  \item \textbf{Posture and monitoring} (Tenet~5): endpoint posture signals and the observability layer continuously measure asset integrity.
  \item \textbf{Telemetry to improve posture} (Tenet~7): an event pipeline collects decision and access telemetry used to refine policy.
\end{enumerate}
\noindent The architecture therefore adopts Zero Trust \emph{by SP~800-207 tenet}, while declining the SASE/ZTNA \emph{deployment model} whose supplier-operated data plane is incompatible with the institution's sovereignty objectives. SASE may still be adopted tactically where sovereignty is not at stake; it is not the structural foundation.
\end{answerbox}

% ============================================================================
\section{Questions on governance models}
\label{sec:q-governance}
% ============================================================================

\epigraph{``The price of freedom is eternal vigilance.''}{--- Attributed to Thomas Jefferson}

% --- Q2.1 ---
\begin{questionbox}[Q2.1 --- ``Federated governance creates accountability gaps --- who is responsible when a breach occurs?'']
If departments manage their own zones, the security function cannot be held responsible for their failures. Accountability becomes diffuse and no one is truly in charge.
\end{questionbox}

\begin{answerbox}[Answer to Q2.1]
\textbf{This is a genuine concern} and the dossier addresses it explicitly through \keyword{enclave charters} modelled on Oxford~\cite{oxford-infosec} and ETH~Z\"urich~\cite{eth-baseline}:

\begin{enumerate}[nosep]
    \item Each zone has an \emph{identified owner} (unit head or delegate);
    \item The charter specifies incident contacts, logging baselines, and minimum security standards;
    \item The central security function defines the \emph{standards} and operates the \gls{soc}; departments \emph{implement} within their perimeter;
    \item This is not \emph{less} accountability but \emph{more granular} accountability --- a model where one officer is nominally responsible for everything but cannot actually control departmental behaviour creates the \emph{real} accountability gap.
\end{enumerate}

\noindent At Oxford, unit heads are explicitly \emph{accountable} for information security in their area~\cite{oxford-infosec}. At Michigan, an IT Council coordinates across nineteen schools with clear responsibility lines~\cite{umich-it-council}. At ETH~Z\"urich, the Administrator Rights Agreement~\cite{eth-admin-agreement} is a binding \emph{transfer} of responsibility, not an abdication.

\medskip\noindent\textbf{Key point:} the question is not ``centralised vs.\ distributed accountability'' but ``formal accountability vs.\ informal unaccountability.'' The situation in which users run shadow IT with zero accountability is the real governance failure.

\medskip\noindent\textbf{Live-incident containment authority:} accountability includes the power to act during an active breach. The enclave charter pre-grants the central security function a scoped, auditable, time-bounded \emph{emergency containment authority} --- the right to force network isolation or quarantine of a zone on defined incident criteria, logged and reviewed post-hoc by the digital-policy commission --- overriding zone autonomy in real time. This is not an erosion of the federated model but its safeguard: \gls{nis2} Art.~21 (technical and operational risk-management measures) and Art.~23 (incident-reporting duty) already vest a central duty to contain; NIST CSF~2.0 Respond and ISO/IEC~27035 standardise the containment workflow; and the Oxford and ETH~Z\"urich unit-head models already subordinate the unit to a central isolation decision~\cite{oxford-infosec,eth-baseline}. The authority is operationalised in the Wave~7 audit-and-incident-response runbooks of the Technical Reference, not merely asserted here.
\end{answerbox}

\bigskip

% --- Q2.2 ---
\begin{questionbox}[Q2.2 --- ``Departments will treat autonomy as the absence of responsibility'']
Giving departments autonomy over their IT will lead to neglect: unpatched servers, exposed databases, abandoned research machines connected to the network.
\end{questionbox}

\begin{answerbox}[Answer to Q2.2]
\textbf{This risk is real} and is identified in the dossier's risk register (rated Medium probability / High impact).

\medskip\noindent\textbf{Mitigations are built into the proposal:}
\begin{itemize}[nosep]
    \item \textbf{Charters with accountability:} mandatory incident contacts, logging baseline, periodic reviews;
    \item \textbf{Baseline monitoring by the \gls{soc}:} the central security team monitors all zones for anomalies, including research zones;
    \item \textbf{Default-deny at internal boundaries:} a neglected zone cannot propagate to other zones;
    \item \textbf{ETH~Z\"urich isolation model:} non-compliant resources are placed in isolated network zones~\cite{eth-baseline} so that weaknesses are not exploitable from outside;
    \item \textbf{Periodic audits:} a digital-policy governance body reviews zone compliance.
\end{itemize}

\noindent\textbf{Counter-question:} under the current monolithic model, are departmental ``servers under desks'' patched, monitored, and accountable? The answer is typically \emph{no} --- they exist as shadow IT with \emph{zero} oversight. Formalised autonomy with charters is strictly \emph{better} than informal unmanaged autonomy.
\end{answerbox}

\bigskip

% --- Q2.3 ---
\begin{clarificationbox}[Q2.3 --- ``This proposal undermines the security function's authority'']
By proposing that academic representatives sit on a digital-policy commission, and that an independent commission can override IT policy decisions, the proposal effectively strips the security function of the authority needed to protect the institution.
\end{clarificationbox}

\begin{answerbox}[Answer to Q2.3]
This question rests on a distinction that governance frameworks draw explicitly. The proposal does not remove the security function's authority --- it \emph{positions it correctly} within a governance structure. There is a fundamental distinction between:

\begin{description}[nosep]
    \item[Governance:] strategic direction and risk appetite (the digital-policy commission);
    \item[Management:] operational execution (the security function and IT department).
\end{description}

This distinction is foundational in \gls{cobit}~2019~\cite{cobit} and in the work of Weill and Ross~\cite{weill-ross} on IT governance. De~Haes and Van~Grembergen~\cite{dehaes2015} explicitly warn against conflating governance with management: the security function \emph{executes} security policy; the governance body \emph{defines} the risk appetite and resolves conflicts between security and the institutional mission.

\medskip\noindent No security officer at MIT, Stanford, Oxford, Michigan, or ETH~Z\"urich has \emph{unilateral} authority over all IT decisions~\cite{umich-it-council, oxford-infosec}. In all these institutions, governance bodies with academic representation define policy, and the security function implements it. This is not a diminishment --- it is professional governance.

\medskip\noindent\textbf{The real question:} should a single operational officer have the authority to make decisions affecting academic freedom, research competitiveness, and the fundamental rights of employees \emph{without external oversight}? Every governance framework says no.
\end{answerbox}

\bigskip

% --- Q2.4 ---
\begin{clarificationbox}[Q2.4 --- ``Academics should not interfere with security --- they lack expertise'']
Security is a technical discipline. Putting academics on a digital-policy commission is like putting passengers in the cockpit.
\end{clarificationbox}

\begin{answerbox}[Answer to Q2.4]
A more precise analogy might be surgeons having input on hospital operating-theatre design --- academics are \emph{the people whose work the infrastructure exists to serve}, not passive passengers. No hospital would exclude surgeons from that process and defer entirely to the facilities manager.

\medskip\noindent\textbf{The expertise argument conflates two levels:}
\begin{itemize}[nosep]
    \item \emph{Technical implementation} (firewall rules, \gls{ids} configuration): this is the security function's domain --- no one disputes this;
    \item \emph{Policy direction} (what should be restricted, what autonomy is needed, what risk appetite is acceptable): this is a \emph{governance} question requiring input from all stakeholders.
\end{itemize}

\noindent NIST CSF~2.0 Govern (GV) function~\cite{nist-csf} explicitly integrates ``stakeholder needs'' into cybersecurity strategy. ISO~27001 Clause~4.2~\cite{iso27001} requires understanding ``the needs and expectations of interested parties.'' EDUCAUSE~\cite{educause-governance} recommends governance structures including academic representation. Excluding the primary users of the infrastructure from governance decisions violates every major IT governance framework.

\medskip\noindent Furthermore, in a university context, many academics \emph{are} technical experts: computer-science researchers, cybersecurity specialists, network engineers. Dismissing their expertise because they are ``academics'' is a category error.
\end{answerbox}

% ============================================================================
\section{Questions on research zone autonomy and local administrative privileges}
\label{sec:q-autonomy}
% ============================================================================

\epigraph{``Discovery consists in seeing what everybody else has seen and thinking what nobody has thought.''}{--- Attributed to Erwin Schr\"odinger}

% --- Q3.1 ---
\begin{questionbox}[Q3.1 --- ``Giving researchers root access creates unacceptable risk'']
Root access means researchers can disable security tools, install vulnerable software, and create backdoors. One compromised research machine can become a launchpad for attacks against the entire institution.
\end{questionbox}

\begin{answerbox}[Answer to Q3.1]
\textbf{Partial merit:} root access is indeed a powerful privilege that can be misused. This is why the dossier does not propose \emph{unrestricted} root access but \emph{accountable} root access within dedicated zones.

\medskip\noindent\textbf{The ETH~Z\"urich model~\cite{eth-admin-agreement}:}
\begin{enumerate}[nosep]
    \item The user requests root rights for a documented purpose (thesis, research project);
    \item By signing the agreement, the user accepts \emph{personal responsibility} for maintenance;
    \item The user commits not to disable basic security mechanisms (logs, updates);
    \item Non-compliant machines are placed in isolated network zones~\cite{eth-baseline};
    \item It is a \emph{transfer of responsibility}, not an abdication of control.
\end{enumerate}

\noindent\textbf{Why this is safer than the alternative:} when root access is denied, researchers jailbreak managed machines, use hidden VMs, connect lab equipment via personal hotspots, or deploy entirely unmanaged servers under desks. These shadow systems have \emph{no} logging, \emph{no} patching, \emph{no} accountability, and \emph{no} isolation. Formal root access with charters, accountability, and network isolation is strictly safer than informal root access with no controls.

\medskip\noindent\textbf{The ``launchpad'' fear:} in a segmented architecture, a compromised research machine cannot reach the administrative enclave --- this is the entire point of segmentation. At CERN, experiment networks are isolated from the general network by strict firewalls and \gls{acl}s~\cite{cern-tn}, yet physicists retain extensive local control.
\end{answerbox}

\bigskip

% --- Q3.2 ---
\begin{questionbox}[Q3.2 --- ``Researcher autonomy = researcher liability --- this unfairly burdens academics'']
By transferring responsibility via charters, the institution offloads security risk onto academics who are not trained for it. This is unfair and possibly illegal under employment law.
\end{questionbox}

\begin{answerbox}[Answer to Q3.2]
\textbf{This concern has merit when pushed to extremes} --- researchers should not be expected to be full-time system administrators. The dossier addresses this through multiple mechanisms:

\begin{itemize}[nosep]
    \item Central IT still provides the backbone, \gls{iam}, \gls{soc}, and baseline monitoring --- researchers are not ``on their own'';
    \item ETH~Z\"urich deploys \emph{Data Stewards}~\cite{eth-data-stewardship} and departmental IT staff (ISGs)~\cite{eth-isg-infk} who bridge the gap between central IT and researchers;
    \item The charter defines \emph{proportionate} responsibilities: a researcher who installs a notebook server has lighter obligations than one running a public-facing web service;
    \item The progressive transition includes training of local IT correspondents by the IT services.
\end{itemize}

\noindent\textbf{On the legal point:} the charter model is an established, documented practice at institutions such as ETH~Z\"urich and Oxford~\cite{eth-admin-agreement, oxford-infosec}. The responsibility is \emph{contractual and proportionate}, not a blanket transfer of institutional liability.

\medskip\noindent\textbf{Counter-question:} under the current model, when a researcher deploys shadow IT that is breached, who is liable? The answer is unclear --- which is precisely the governance failure that charters resolve.
\end{answerbox}

\bigskip

% --- Q3.3 ---
\begin{questionbox}[Q3.3 --- ``The Science DMZ bypasses security controls --- this is a security hole by design'']
The \gls{sciencedmz} explicitly removes the enterprise firewall from the research data path. This means research traffic is not inspected. This is an indefensible gap in security posture.
\end{questionbox}

\begin{answerbox}[Answer to Q3.3]
\textbf{The premise is factually incorrect.} The \gls{sciencedmz} does \emph{not} remove security --- it replaces one type of security (a stateful firewall, designed for enterprise traffic) with another appropriate for research traffic:

\begin{description}[nosep]
    \item[\gls{acl}s on routers:] whitelist-based filtering at line speed;
    \item[Zeek (formerly Bro) \gls{ids}:] a dedicated intrusion-detection system monitoring all \gls{sciencedmz} traffic~\cite{paxson1999};
    \item[perfSONAR:] active performance monitoring that detects anomalies~\cite{perfsonar};
    \item[\gls{dtn} hardening:] Data Transfer Nodes are purpose-built, minimal-service machines with a reduced attack surface.
\end{description}

This model was developed by ESnet~\cite{dart2013}, the U.S.\ Department of Energy's network serving nuclear laboratories. It is deployed at \textbf{over 100 universities} through NSF funding~\cite{esnet-history,nsf-cc}. If bypassing the enterprise firewall were a ``security hole by design,'' neither the DoE nor the NSF --- both security-conscious organisations --- would have funded and promoted it.

\medskip\noindent\textbf{The deeper point:} a stateful firewall that, in a documented case, mishandled TCP window scaling on high-volume long-distance flows and degraded research throughput by \textbf{five- to twelve-fold}~\cite{esnet-pennstate} does not \emph{provide} security. It provides the \emph{appearance} of security while driving researchers to bypass it via shadow IT. A \gls{sciencedmz} with \gls{acl}s, \gls{ids}, and monitoring provides \emph{actual} security appropriate to the traffic profile.
\end{answerbox}

\bigskip

% --- Q3.4 ---
\begin{questionbox}[Q3.4 --- ``If the Science DMZ were so good, our network operator would already have deployed it'']
The model is presented as state of the art and adopted by 100-plus universities. Yet neither our national research network nor our HPC team publishes any Science DMZ today. The absence is itself an argument: if it were the right answer for us, our \gls{nren} or our \gls{hpc} team would already have built it.
\end{questionbox}

\begin{answerbox}[Answer to Q3.4]
\textbf{The factual premise may be correct; the inference is not --- this is an appeal to absence.} That an institution or its \gls{nren} has not yet deployed a pattern is not evidence that the pattern is unsuitable; it is, in most cases, evidence about \emph{capacity and funding routes}.

\medskip\noindent Three observations frame the diagnostic generically:
\begin{enumerate}[nosep]
    \item \textbf{The substrate is often already present.} High-bandwidth \gls{nren} access (e.g.\ via \gls{geant}~\cite{geant}), a national \gls{csirt} for incident response, and an \gls{hpc} facility moving large volumes of data internationally are frequently in place. What is missing is the architectural \emph{pattern}, not the bandwidth. Current research transfer over generic enterprise paths (\texttt{rsync}, \texttt{scp}) is an enterprise-IT profile, not a research-tuned one~\cite{dart2013, esnet-cases}.
    \item \textbf{Non-deployment tracks capacity, not merit.} A small research ecosystem with a limited pool of research-network engineers, and the absence of a dedicated funding programme equivalent to the NSF CC* line that explicitly funded Science DMZs at over 100 U.S.\ campuses~\cite{esnet-history,nsf-cc}, explains non-deployment without implying technical inadequacy.
    \item \textbf{Convergent practice elsewhere is overwhelming.} ESnet, NSF, \gls{geant}, and the Medical Science DMZ pattern --- documented in implemented deployments at major U.S.\ biomedical-research institutions~\cite{medical-sdmz} --- converge on the design. ``We have not done X yet'' is a weak argument against X when every comparable research ecosystem has.
\end{enumerate}

\medskip\noindent\textbf{The realistic form is mutualised.} The proposal is not a siloed institution-only Science DMZ but a federated research-network zone --- with a Medical Science DMZ enclave for biomedical workloads --- \emph{co-designed} across the research-active institutions and the \gls{nren}. Mutualisation is the form in which the operational and economic case closes at small scale.

\medskip\noindent\textbf{One honest caveat.} A 2024 quantitative study confirms that Science~DMZ benefits are not automatic: without deliberate tuning of hosts, protocol stacks and measurement, latency on long-RTT paths can even degrade~\cite{sciencedmz-perf}. This argues for operational ownership and tuning, not for inaction.
\end{answerbox}

% ============================================================================
\section{Questions on BYOD and AI cloud services}
\label{sec:q-byod-ai}
% ============================================================================

\epigraph{``The best way to predict the future is to invent it.''}{--- Alan Kay, Xerox PARC, 1971}

% --- Q4.1 ---
\begin{questionbox}[Q4.1 --- ``BYOD is inherently insecure --- personal devices cannot be trusted'']
Personal devices are not managed, may be infected, may lack encryption, and may be shared with family members. Allowing them on the network is an unnecessary risk.
\end{questionbox}

\begin{answerbox}[Answer to Q4.1]
\textbf{\gls{byod} is not a policy choice --- it is a reality.} Faculty members use personal laptops. Researchers use personal phones. Visiting scholars arrive with their own devices. The question is not whether \gls{byod} exists but whether it is \emph{managed or unmanaged}.

\medskip\noindent The dossier proposes \emph{progressive interoperability}:
\begin{enumerate}[nosep]
    \item \textbf{Without authentication:} Internet only (guest zone);
    \item \textbf{With \gls{mfa} authentication:} access to teaching and research resources matching the user's profile;
    \item \textbf{With a compliant device (verified posture):} extended access including administrative resources via a secure \gls{vdi}.
\end{enumerate}

This is precisely the \gls{zt} approach recommended by NIST SP~800-207~\cite{nist-zt}: access decisions based on identity and device posture, not on whether the device is ``managed'' or ``personal.''

\medskip\noindent\textbf{\gls{eduroam}}~\cite{eduroam} --- the global academic federated Wi-Fi system --- has operated successfully for two decades on exactly this principle: authenticate users, not devices.

\medskip\noindent The alternative --- banning \gls{byod} --- pushes devices onto personal hotspots and creates genuinely unmanaged, invisible connections. This is \emph{less} secure than a dedicated \gls{byod} subnet with progressive access.
\end{answerbox}

\bigskip

% --- Q4.2 ---
\begin{questionbox}[Q4.2 --- ``Giving researchers free access to non-EU AI cloud services creates GDPR liability'']
Commercial AI cloud services are not fully \gls{gdpr}-compliant. If researchers send personal data to these services, the institution faces regulatory sanctions.
\end{questionbox}

\begin{answerbox}[Answer to Q4.2]
\textbf{This is a legitimate concern} and the dossier does not propose unfettered access without safeguards. It proposes a \emph{differentiated framework}:

\begin{enumerate}[nosep]
    \item \textbf{Data classification:} personal data is \emph{never} sent to a third-country service without a lawful basis and a valid transfer mechanism --- an adequacy route where one exists (e.g.\ the EU--US Data Privacy Framework for certified providers), or standard contractual clauses (\gls{scc}s) supplemented, per \emph{Schrems~II}, by a documented Transfer Impact Assessment (\gls{tia}) and supplementary technical measures (encryption, pseudonymisation, key custody), as operationalised in the regulatory mapping;
    \item \textbf{Environment separation:} research AI environments are isolated from the administrative zone, preventing accidental mixing of personal data with research API calls;
    \item \textbf{\gls{aiact} risk-based compliance:} for researchers \emph{using} commercial AI services, the applicable framework is the \gls{aiact}'s risk-based classification --- \emph{not} a blanket research exemption; the Art.~2(6)/2(8) exemptions cover bespoke, pre-market research tools, not commercial-API use~\cite{eu-ai-act};
    \item \textbf{Traceability:} logging of external API access for auditing;
    \item \textbf{Training:} researcher awareness of \gls{gdpr} implications.
\end{enumerate}

\noindent\textbf{The alternative is worse:} when access is blocked, researchers use personal accounts on personal devices outside the institutional network, sending the same data with \emph{zero} institutional safeguards, \emph{zero} \gls{scc}s, and \emph{zero} traceability. The institution's \gls{gdpr} exposure is \emph{higher}, not lower, when it refuses to provide a managed pathway.

\medskip\noindent\textbf{Competitive reality:} leading research universities increasingly provide researchers with \emph{managed} access to global AI platforms rather than prohibiting them outright. An institution that blocks such access risks putting its researchers at a competitive disadvantage relative to peers who provide a governed pathway.
\end{answerbox}

\bigskip

% --- Q4.3 ---
\begin{clarificationbox}[Q4.3 --- ``The AI access proposal is about convenience, not necessity'']
Researchers managed perfectly well before generative AI. AI cloud access is a luxury, not a structural need for research. The institution should not compromise its security posture for fashionable tools.
\end{clarificationbox}

\begin{answerbox}[Answer to Q4.3]
This perspective may underestimate the pace of change in research methodologies. Behavioural research on decision-making~\cite{samuelson-zeckhauser} documents \emph{status-quo bias}: the tendency to prefer the current state of affairs simply because it is familiar, regardless of evolving evidence.

\medskip\noindent\textbf{On the facts:} foundation models and generative AI are \emph{not} optional luxuries for research in computer science, NLP, machine learning, data science, computational biology, digital humanities, and increasingly the social sciences. Major conferences (NeurIPS, ICML, ACL, CVPR) increasingly feature work relying on cloud AI platforms, and funding agencies explicitly fund AI-based research.

\medskip\noindent Dismissing AI as ``fashionable tools'' is analogous to dismissing the Internet in 1995 or statistical software in 1985. The question is not \emph{whether} universities will use AI cloud services but \emph{whether they will do so through managed, secure channels or through ungoverned shadow IT}.

\medskip\noindent\textbf{Structural need:} AI research \emph{requires} access to APIs (for model comparison, fine-tuning, evaluation), cloud accelerators (for training), and open-source model repositories. These are not interchangeable with local-only tools. Teaching advanced AI courses without API access is like teaching chemistry without a laboratory. The differentiating factor is not the tool but the \emph{governed workflow}: managed access with classification and logging improves outcomes, while a blanket ban displaces the same activity into the ungoverned shadow.
\end{answerbox}

% ============================================================================
\section{Questions on institutional comparisons and transferability}
\label{sec:q-comparisons}
% ============================================================================

\epigraph{``In theory, there is no difference between theory and practice. In practice, there is.''}{--- Attributed to Yogi Berra}

% --- Q5.1 ---
\begin{clarificationbox}[Q5.1 --- ``You are comparing us to MIT and Stanford --- this is absurd'']
MIT has a multi-billion endowment; Stanford has thousands of IT staff. We are a small institution. These comparisons are meaningless and the models cannot transfer.
\end{clarificationbox}

\begin{answerbox}[Answer to Q5.1]
The comparison is not about replicating budgets but about adopting scale-independent architectural principles. The question assumes that the only relevant difference is scale, and that smaller institutions cannot adopt principles from larger ones --- a \emph{false equivalence} between ``adopting a principle'' and ``replicating a budget''.

\medskip\noindent\textbf{The dossier does not propose replicating any institution's budget.} It proposes adopting \emph{architectural principles} (segmentation, federation, charters) that are \emph{scale-independent}:
\begin{itemize}[nosep]
    \item \textbf{\gls{vlan} segmentation} requires no additional \emph{hardware} budget --- it is a configuration decision on existing switches; the cost is design and operational labour, not procurement;
    \item \textbf{Charters} are documents, not infrastructure purchases;
    \item \textbf{The \gls{sciencedmz}} uses open-source tools (perfSONAR, Zeek) and \gls{nren} services already available through \gls{geant}~\cite{geant, perfsonar};
    \item \textbf{Federated governance} is an organisational choice, not a technology procurement.
\end{itemize}

\noindent\textbf{The evidence base deliberately spans institutions of different sizes:} not only MIT and Stanford, but also smaller European universities, ETH~Z\"urich, Michigan, NC~State, Penn State and Wisconsin-Madison~\cite{ncsu-governance, pennstate-it, uw-madison, umich-it-council}. The architectural principles are consistent across all sizes.

\medskip\noindent\textbf{The scale argument actually favours federation:} a small institution has \emph{fewer} resources for a monolithic central team to manage everything. Federation distributes the workload to where domain expertise exists.

\medskip\noindent\textbf{The deeper point: these are not larger institutions' choices to borrow, but a minimal necessary set.} The functions the architecture discharges --- authenticating users, protecting sensitive and personal data, segmenting, auditing --- are required of \emph{any} institution offering a digital infrastructure. Several are legally compelled with no size threshold: a data-protection duty applies to any institution processing personal data, irrespective of size, under \loc{data-protection-regime} (in the EU, the \gls{gdpr}~\cite{gdpr}, which applies irrespective of size). The choice a small institution faces is therefore not ``afford the leading universities' model or not'', but ``discharge these already-owed obligations coherently and sovereignly, or implicitly and dependently''. Only the instantiation scales with size; the obligations do not~\cite{ng-tech-ref-2026}.
\end{answerbox}

\bigskip

% --- Q5.2 ---
\begin{questionbox}[Q5.2 --- ``Cherry-picking successful models ignores failed federations'']
The dossier only cites universities where federation works. What about universities that tried federation and failed, reverting to centralisation?
\end{questionbox}

\begin{answerbox}[Answer to Q5.2]
\textbf{This is a fair challenge.} Any honest analysis must consider counter-examples.

\medskip\noindent\textbf{However:}
\begin{enumerate}[nosep]
    \item The evidence base spans \textbf{a broad set of reference universities} across the US, UK, Switzerland and continental Europe, plus CERN, plus the wider \gls{nren} ecosystem. This is not cherry-picking --- it is the \emph{dominant pattern} in high-performing universities worldwide.
    \item The EDUCAUSE Top~10 IT Issues~\cite{educause2024} and EDUCAUSE IT-governance resources~\cite{educause-governance} consistently recommend federated governance for higher education. This is sector consensus, not anecdote.
    \item Counter-examples of ``failed federation'' typically involve institutions that \emph{federated without accountability mechanisms} (no charters, no \gls{soc} oversight, no minimum standards). The charter-based model addresses this risk directly.
\end{enumerate}

\noindent\textbf{The burden of evidence is symmetrical:} if the questioner claims that federation fails, they should cite specific examples and show that the failures were caused by \emph{federation itself} rather than by the absence of the governance mechanisms proposed here.
\end{answerbox}

\bigskip

% --- Q5.3 ---
\begin{clarificationbox}[Q5.3 --- ``Those universities were built federated --- we cannot retrofit federation onto a centralised institution'']
MIT, Cambridge, and Stanford were \emph{born} decentralised. Their federation is historical, not deliberate. You cannot impose federation on an institution that has been centralised for years.
\end{clarificationbox}

\begin{answerbox}[Answer to Q5.3]
\textbf{Factually incorrect for several documented universities:}
\begin{itemize}[nosep]
    \item \textbf{NC~State} explicitly \emph{redesigned} its IT governance for digital transformation, as documented in EDUCAUSE Review~\cite{ncsu-governance};
    \item \textbf{Penn State} \emph{unified} IT professionals from centralised units while introducing participatory governance and SLAs~\cite{pennstate-it};
    \item \textbf{Wisconsin-Madison} adopted federation as a deliberate ``modernisation and federation'' strategy~\cite{uw-madison};
    \item \textbf{University of Michigan} evolved from a more centralised model to formal federation with the creation of an IT Council~\cite{umich-it-council}.
\end{itemize}

\noindent Federation is not a ``natural state'' --- it is a \emph{governance choice} that can be adopted, with appropriate transition planning. The progressive transition plan (governance first, then pilot on selected entities, then gradual extension) is designed precisely for institutions moving \emph{from} centralisation \emph{to} federation.
\end{answerbox}

\bigskip

% --- Q5.4 ---
\begin{questionbox}[Q5.4 --- ``The reversibility claim rests on soft, single-vendor conformance tests'']
The lock-in answer rests on exit tests and a reversibility playbook; if those tests are anchored on single-vendor specifications, the guarantee is soft.
\end{questionbox}

\begin{answerbox}[Answer to Q5.4]
\textbf{The exit tests are cross-implementation by design, not single-vendor.} Each replaceable layer is tested against an \emph{alternative} implementation: \gls{scim} user/group export to an alternative identity provider, policy bundles re-imported into an alternative engine, container images to an alternative OCI registry, switch configuration to an alternative-vendor device, the analytical store re-pointed, and object-store operations against an independent S3-compatible implementation. Reversibility is therefore tested behaviour, not a vendor promise --- developed in full in the standalone Technical Reference~\cite{ng-tech-ref-2026}.
\end{answerbox}

% ============================================================================
\section{Questions on the independent adjustments commission}
\label{sec:q-iac}
% ============================================================================

\epigraph{``Injustice anywhere is a threat to justice everywhere.''}{--- Martin Luther King Jr., \textit{Letter from Birmingham Jail}, 1963}

% --- Q6.1 ---
\begin{questionbox}[Q6.1 --- ``An independent adjustments commission creates bureaucratic overhead and slows down IT decisions'']
Adding an independent commission that can override IT policy decisions introduces delays and bureaucracy into what should be agile operational decisions.
\end{questionbox}

\begin{answerbox}[Answer to Q6.1]
\textbf{The \gls{icradp} does not ``override'' operational decisions.} It provides an \emph{adjustment} mechanism for \emph{individual} cases where a \emph{general} policy creates a disproportionate obstacle for a specific person. This is analogous to a court system: the existence of courts does not slow down legislation; it corrects general rules that produce unjust outcomes in specific cases.

\medskip\noindent\textbf{On timing:}
\begin{itemize}[nosep]
    \item The commission rules within a short, published deadline --- faster than many IT change-management processes;
    \item In documented emergencies (medical, professional), a provisional adjustment can be granted within a sharply shorter window;
    \item The commission handles \emph{exceptions}, not routine operations. The segmented architecture, by providing appropriate zones for most needs, \emph{reduces} the number of adjustment requests.
\end{itemize}

\noindent\textbf{Cost of the alternative:} without a formal mechanism, adjustment requests are handled through \emph{informal negotiations} with the IT department --- a process that is slower, less predictable, less fair, and produces no auditable record. The \gls{icradp} \emph{streamlines} exception handling by centralising and formalising it.
\end{answerbox}

\bigskip

% --- Q6.2 ---
\begin{clarificationbox}[Q6.2 --- ``The commission is a Trojan horse to override IT security decisions'']
Under the guise of ``reasonable adjustments,'' every dissatisfied user will appeal to the commission to bypass security controls. The commission will become a mechanism for institutional circumvention of IT policy.
\end{clarificationbox}

\begin{answerbox}[Answer to Q6.2]
This concern is understandable; the following structural safeguards are designed precisely to prevent abuse:

\begin{enumerate}[nosep]
    \item \emph{Documented justification:} requests must specify medical, personal, professional, or fundamental-rights grounds;
    \item \emph{Proportionality assessment:} the commission evaluates whether the adjustment creates a disproportionate security risk (not merely an inconvenience for IT);
    \item \emph{Reasoned refusals:} the commission can and will refuse requests that do not meet the criteria;
    \item \emph{Annual transparency report:} statistics on granted/refused requests, preventing either systematic overriding or systematic refusal;
    \item \emph{The security function is heard:} the security perspective is formally considered in every decision.
\end{enumerate}

\noindent\textbf{Precedent:} reasonable-adjustment committees that support students and staff operate on similar principles across many institutions and have not devolved into mechanisms for ``bypassing'' academic standards. Courts, ombudsmen, and appeal mechanisms exist in every well-governed institution without destroying the authority of the decision-makers they review.

\medskip\noindent\textbf{The deeper issue:} an IT department that cannot withstand \emph{reasoned, independent review} of its policies is an IT department whose policies may not be defensible. Accountability mechanisms strengthen governance; they do not undermine it.
\end{answerbox}

\bigskip

% --- Q6.3 ---
\begin{clarificationbox}[Q6.3 --- ``Citing fundamental rights for IT policy is emotional manipulation'']
Invoking the EU Charter of Fundamental Rights, the ECHR, and human-rights case law for a discussion about firewall rules is disproportionate and emotionally manipulative. This conflates totalitarian censorship with ordinary IT administration.
\end{clarificationbox}

\begin{answerbox}[Answer to Q6.3]
The legal framework cited is binding, not rhetorical. Here is how it applies concretely:
\begin{itemize}[nosep]
    \item In EU jurisdictions, Article~13 of the EU Charter protects academic freedom as a \emph{fundamental right} and applies to \emph{universities} --- that is not a stretch; it is the literal scope of the provision. Comparable jurisdictions protect academic freedom under their own constitutional or statutory instruments.
    \item The \gls{ecthr} in \emph{Sorgu\c{c} v.\ Turkey}~\cite{sorguc2009} and \emph{Erdo\u{g}an v.\ Turkey}~\cite{erdogan2014} applied freedom-of-expression protections specifically to \emph{academic} contexts.
    \item The \gls{cjeu} in \emph{Scarlet Extended v.\ SABAM}~\cite{scarlet2011} ruled that \emph{systematic filtering} of Internet traffic is incompatible with Article~11 of the Charter.
    \item The Grand Chamber in \emph{B\u{a}rbulescu v.\ Romania}~\cite{barbulescu2017} established a \textbf{six-criteria proportionality test} for digital surveillance \emph{in the workplace} --- the exact context of university IT policies.
\end{itemize}

\noindent These are not abstract philosophical references. They are \emph{binding law} in the EU and Council of Europe member states; where the adopter is so bound, an IT policy that systematically restricts digital access and monitors communications \emph{must} satisfy the proportionality requirements of these instruments. Elsewhere, they stand as a strongly persuasive proportionality benchmark alongside the adopter's own fundamental-rights instruments.

\medskip\noindent\textbf{Scale of the claim:} the dossier does not compare university IT policies to totalitarian censorship. It argues that digital policies \emph{engage} fundamental rights and must therefore pass a proportionality test. This is a measured, legally grounded claim --- not emotional manipulation.
\end{answerbox}

\bigskip

% --- Q6.4 ---
\begin{questionbox}[Q6.4 --- ``Technostress and moral injury are not real occupational hazards --- this is pseudoscience'']
Technostress is a buzzword. Moral injury applies to soldiers, not to academics who cannot install software. Citing these concepts to argue against IT policies is scientifically dubious.
\end{questionbox}

\begin{answerbox}[Answer to Q6.4]
\textbf{The scientific evidence is extensive and peer-reviewed:}

\begin{description}[nosep]
    \item[Technostress:] La~Torre et al.~\cite{latorre2019} published a systematic review of \textbf{105 studies} in the \emph{International Archives of Occupational and Environmental Health} (Springer), establishing technostress as a documented occupational risk with measurable effects on well-being and productivity. Tarafdar et al.~\cite{tarafdar2007} demonstrated its impact on role stress and performance in the \emph{Journal of Management Information Systems}. Ragu-Nathan et al.~\cite{ragunathan2008} validated the construct empirically in \emph{Information Systems Research}. These are top-tier peer-reviewed journals, not ``buzzword'' sources.

    \item[Moral injury:] Litz et al.~\cite{litz2009} introduced the construct in the \emph{Clinical Psychology Review}. Williamson et al.~\cite{williamson2018} published a systematic review and meta-analysis in the \emph{British Journal of Psychiatry}, extending the concept to \emph{occupational} contexts. Hanna et al.~\cite{hanna2023} specifically studied moral injury in \emph{university staff}. The concept is not limited to military contexts --- it is a well-established framework in occupational psychology.
\end{description}

\noindent Dismissing peer-reviewed systematic reviews as ``pseudoscience'' without engaging the evidence is itself an unscientific position. One may disagree with the \emph{policy implications} drawn from this literature, but denying the documented phenomena requires counter-evidence, not dismissal. These findings describe documented risks; the dossier pairs them with the mitigation it proposes (proportionate policy and an adjustments mechanism), not with a claim that restriction is always harmful.
\end{answerbox}

% ============================================================================
\section{Questions prompted by recent incidents}
\label{sec:q-incidents}
% ============================================================================

\epigraph{``Data is not information, information is not knowledge, knowledge is not understanding, understanding is not wisdom.''}{--- Russell L.\ Ackoff, \textit{From Data to Wisdom}, 1989}

% --- Q7.1 ---
\begin{clarificationbox}[Q7.1 --- ``Recent university breaches prove we need MORE restriction, not less'']
The headlines are clear: universities are under siege from ransomware and supply-chain compromises. This is not the time to relax security by giving departments autonomy.
\end{clarificationbox}

\begin{answerbox}[Answer to Q7.1]
This is a frequently heard interpretation that warrants careful factual examination. Sector incident analysis points consistently to a small set of root causes:

\begin{itemize}[nosep]
    \item \textbf{Third-party software vulnerabilities (supply chain):} neither ``more restriction'' on research zones nor ``less department autonomy'' would have prevented them. The mitigation is vendor risk management, patching, and \emph{isolation of critical applications} --- which is what segmentation achieves.
    \item \textbf{Centralised administrative perimeters:} ransomware that exploits a single undifferentiated perimeter is amplified by the \emph{absence} of segmentation, where email, data servers, and administrative services share one perimeter.
    \item \textbf{Credential compromise and lateral movement:} segmentation is the standard defence against lateral movement~\cite{verizon-dbir2024}.
    \item \textbf{Contained incidents:} where an incident \emph{is} confined to a subnet, that is a demonstration that segmentation \emph{works}~\cite{rennes-cyber}.
\end{itemize}

\noindent\textbf{None of these incident classes were caused by ``openness,'' ``federation,'' or ``research-zone autonomy.''} They were caused by third-party vulnerabilities, credential compromise without segmentation, and centralised perimeters that became single points of failure. The Sophos State of Ransomware in Education 2025~\cite{sophos-ransomware-edu2025} identifies ``exploited vulnerabilities'' and ``unknown security gaps'' as top causes --- consistent with patching and segmentation, not with restriction. The proposed mechanisms (segmentation, enclaves, charters, zone-specific monitoring) \emph{directly address} every one of these root causes; the ``more restriction'' response does not.
\end{answerbox}

\bigskip

% --- Q7.2 ---
\begin{questionbox}[Q7.2 --- ``The incident analysis is biased --- it selectively interprets failures to support its thesis'']
Any incident can be interpreted multiple ways. Claiming that a centralised-perimeter breach proves the need for \emph{less} centralisation is tendentious when it could equally prove the need for \emph{stronger} centralisation.
\end{questionbox}

\begin{answerbox}[Answer to Q7.2]
\textbf{This is a fair methodological challenge.} Let us examine it rigorously.

\medskip\noindent The claim that a centralised-perimeter breach proves ``stronger centralisation'' would need to show that (i) the breach was caused by a \emph{decentralised} part of the infrastructure, and (ii) a more centralised architecture would have prevented the specific attack vector.

\medskip\noindent The available evidence on the canonical cases shows the opposite: such attacks target \emph{central administrative services} (email, data servers), and the damage is amplified because these services sit in a single, undifferentiated perimeter. Centralising \emph{more} services into that perimeter would have \emph{increased} the blast radius.

\medskip\noindent\textbf{The interpretation aligns with independent sector analysis:}
\begin{itemize}[nosep]
    \item Verizon DBIR 2024~\cite{verizon-dbir2024}: credential abuse and ransomware as the dominant measured breach patterns, whose in-network progression segmentation is designed to interrupt (per the NIST/ISO controls above);
    \item NIST SP~800-215~\cite{nist-sp800-215}: a guide to securing enterprise networks through segmentation;
    \item Sophos State of Ransomware in Education 2025~\cite{sophos-ransomware-edu2025}: ``exploited vulnerabilities'' and ``unknown security gaps'' as top causes --- consistent with patching and segmentation, not with restriction.
\end{itemize}

\noindent The analysis is not ``selective interpretation'' --- it is alignment with the independent analysis of Verizon, Sophos, and NIST. The burden is on the questioner to show how ``stronger centralisation'' addresses the documented root causes.
\end{answerbox}

% ============================================================================
\section{Questions on regulatory interpretation}
\label{sec:q-regulatory}
% ============================================================================

\epigraph{``The law is reason, free from passion.''}{--- Aristotle, \textit{Politics}}

% --- Q8.1 ---
\begin{questionbox}[Q8.1 --- ``NIS2 may actually require MORE centralisation, not federation'']
The \gls{nis2} Directive requires incident notification, risk management, and governance. These requirements are easier to demonstrate with a single, centralised IT authority. Federation makes compliance harder.
\end{questionbox}

\begin{answerbox}[Answer to Q8.1]
\textbf{\gls{nis2} is requirements-based, not architecture-prescriptive.} The Directive specifies \emph{what} must be achieved (risk management, incident containment, governance) but does \emph{not} prescribe \emph{how} the infrastructure must be organised.

\medskip\noindent A federated architecture with a central \gls{soc} for detection and notification, per-zone risk assessments and incident playbooks, charters documenting responsibilities and contacts, and a digital-policy commission for strategic oversight, meets every \gls{nis2} requirement while \emph{also} providing the compartmentalisation that \emph{improves} incident containment --- a core \gls{nis2} objective.

\medskip\noindent\textbf{One locatable clock owner.} Where the adopter is an essential or important entity under \loc{cybersecurity-baseline-regime}, the central \gls{soc} owns the incident-reporting clock --- in the EU, the \gls{nis2} Article~23 timeline of a 24-hour early warning, a 72-hour notification and a one-month final report~\cite{nis2} --- while the zone charters oblige each zone to report inward to the \gls{soc} within a sub-deadline. The single-owner pattern is generic: it applies to whatever incident-reporting regime governs the adopter, so federation does not fragment the reporting obligation but gives it a single, accountable owner.

\medskip\noindent\textbf{ISO~27001 precedent:} universities that achieve ISO~27001 certification commonly scope the certification to \emph{central IT services} rather than the entire institution. Clause~4.3 of the standard~\cite{iso27001} expressly permits such a bounded \gls{isms} scope, and certifiers accept this \emph{de facto} federation as a valid scope definition.

\medskip\noindent\textbf{The evidence argument:} a multi-zone architecture with per-zone logging, incident playbooks and regular testing \emph{produces more auditable evidence} than a single monolithic perimeter.
\end{answerbox}

\bigskip

% --- Q8.2 ---
\begin{questionbox}[Q8.2 --- ``The GDPR `risk-based' argument is too convenient --- it could justify anything'']
Saying that \gls{gdpr} Article~32 requires measures ``appropriate to the risk'' does not mean an institution can decide to accept more risk. A supervisory authority could still sanction insufficient security even if the institution claims it assessed the risk.
\end{questionbox}

\begin{answerbox}[Answer to Q8.2]
\textbf{Correct --- the risk-based approach is not a carte blanche.} \gls{gdpr} Article~32 requires measures appropriate to the risk, taking into account the ``state of the art, costs, context, and the likelihood/severity of risks''~\cite{gdpr}. The institution must demonstrate that its choices are \emph{justified}, not merely \emph{convenient}.

\medskip\noindent\textbf{However, this argument applies equally to the monolithic approach:}
\begin{itemize}[nosep]
    \item A monolithic approach that \emph{over-restricts} research and teaching while \emph{under-protecting} the administrative enclave (because the same firewall handles both) may also fail the ``appropriate to the risk'' test --- for different reasons;
    \item The ``state of the art'' element explicitly requires considering what comparable institutions do. The state of the art in higher education \emph{is} segmented, federated architecture; an institution that ignores it may be \emph{more} exposed to regulatory challenge, not less;
    \item The \gls{edpb}~\cite{edpb-guidelines} and \gls{enisa}~\cite{enisa-risk} confirm the risk-based interpretation and that measures must be \emph{adapted to context} --- not uniformly maximal.
\end{itemize}

\noindent The dossier does not argue for ``less security.'' It argues for \emph{differentiated} security: maximum where the risk is highest (administrative data), proportionate where the risk is lower (research). This is precisely what \gls{gdpr} Article~32 contemplates.
\end{answerbox}

\bigskip

% --- Q8.3 ---
\begin{questionbox}[Q8.3 --- ``The proposal never states a GDPR lawful basis or a DPIA trigger'']
The proposal introduces new processing of personal data --- per-zone monitoring, the audit pipeline (which logs subject and decision identifiers), AI-API access logging, and federated identity --- yet it never states the \gls{gdpr} Article~6 lawful basis, nor identifies an Article~35 DPIA trigger. Under the Article~5(2) accountability principle, a controller must be able to demonstrate its lawful basis; as written, a \gls{dpo} could refuse the proposal on the ground that none is stated.
\end{questionbox}

\begin{answerbox}[Answer to Q8.3]
\textbf{The concern is well founded, and the lawful basis is stated here explicitly.} The proposal does introduce new processing, and under \gls{gdpr} Article~5(2) the controller must be able to demonstrate its lawful basis~\cite{gdpr}.

\medskip\noindent\textbf{Lawful basis.} The controller identifies a lawful basis under \loc{data-protection-regime}: a public-mission body typically relies on the public-interest task, a private body on legitimate interests or consent. In the EU instantiation, for a public university the monitoring, audit and identity processing rests on \gls{gdpr} Article~6(1)(e) --- processing necessary for a task carried out in the public interest (the institution's research and teaching mission and its statutory security obligations, including those under \gls{nis2}) --- and, where appropriate, Article~6(1)(f) legitimate interests (network and information security, expressly recognised in Recital~49), subject to the balancing test against data subjects' rights and freedoms. Any processing of special categories of data is conditioned on a corresponding Article~9 ground.

\medskip\noindent\textbf{DPIA.} Because the monitoring and audit pipeline use new technologies and entail systematic processing likely to result in high risk, they are expected to warrant a DPIA under Article~35(1)~\cite{gdpr}. The proposal feeds that assessment rather than replacing it: the per-scenario audit records supply the systematic description, the necessity-and-proportionality analysis and the risk treatment an Article~35 DPIA requires. The per-processing lawful basis is finalised in the DPIA and the Record of Processing Activities (Article~30), owned by the institution's \gls{dpo}.

\medskip\noindent\textbf{This strengthens, rather than weakens, the proposal.} A monolithic perimeter that monitors all traffic uniformly faces the \emph{same} Article~6 and Article~35 obligations with \emph{less} ability to scope and minimise. The segmented, context-bounded design keeps monitoring proportionate --- per-zone scope, minimised audit fields, access-controlled logs --- which is exactly what the necessity-and-proportionality test of Article~35 rewards.
\end{answerbox}

\bigskip

% --- Q8.4 ---
\begin{questionbox}[Q8.4 --- ``Barbulescu is wielded outward but never applied to the proposal's own monitoring'']
Q6.3 invokes the B\u{a}rbulescu six-criteria proportionality test against restrictive monitoring, but the proposal's own per-zone monitoring, \gls{soc} oversight and audit pipeline (logging subject and decision identifiers) are themselves surveillance --- which the proposal never runs through the same six criteria. A reviewer could call this an asymmetry.
\end{questionbox}

\begin{answerbox}[Answer to Q8.4]
\textbf{Fair point --- the same test must, and does, apply to the proposal's own monitoring.} The architecture's monitoring is surveillance in the B\u{a}rbulescu~\cite{barbulescu2017} sense and is assessed here against its six criteria:
\begin{enumerate}[nosep]
  \item \textbf{Prior notification:} monitoring is disclosed in advance through the zone charters and the institution's transparency notices.
  \item \textbf{Extent and intrusion:} the audit pipeline records \emph{metadata} (subject and decision identifiers, verdicts, timestamps), not communication \emph{content}; monitoring is scoped per zone, not blanket, and retained for a bounded period.
  \item \textbf{Legitimate justification:} the aim is network and information security and the institution's \gls{nis2}/\gls{gdpr} duties --- not performance or behavioural surveillance; access to any content would require the weightier justification B\u{a}rbulescu demands.
  \item \textbf{Least intrusive means:} context-bounded, metadata-level monitoring is \emph{less} intrusive than the blanket content inspection a monolithic perimeter tends toward.
  \item \textbf{Consequences and use:} logs feed security operations and incident response under access controls; they are not a routine input to disciplinary or performance decisions.
  \item \textbf{Safeguards:} \gls{soc} access controls, retention limits, \gls{dpo} oversight and the DPIA (Q8.3) provide the safeguards B\u{a}rbulescu requires --- especially the bar on routine content access.
\end{enumerate}
\noindent Far from exempting itself, the proposal's segmented, metadata-level design is what \emph{lets} it satisfy the B\u{a}rbulescu criteria more readily than the uniform monitoring it replaces.

\medskip
\noindent\textbf{Collective dimension.} B\u{a}rbulescu governs the
\emph{individual} proportionality test; in several jurisdictions a
\emph{collective} step precedes it: consultation or co-determination
of \loc{staff-representation-body} before workplace-monitoring tools
are deployed. The Technical Reference makes that consultation a
Wave-2 precondition where the adopter's labour law so requires, and
the \gls{icradp} seat gives staff a standing voice once the
pipeline operates.
\end{answerbox}

\bigskip

% --- Q8.5 ---
\begin{questionbox}[Q8.5 --- ``Federated procurement leaves NIS2 supply-chain security (Art.~21(2)(d)) without an owner'']
The incident answers address supply-chain risk only \emph{after the fact}. \gls{nis2} Article~21(2)(d) imposes a \emph{forward}, standing duty to manage security in the supply chain and supplier relationships. Where zones and faculties procure their own tools, a security officer could refuse the proposal on the ground that no one owns supply-chain assurance \emph{across} autonomous zones.
\end{questionbox}

\begin{answerbox}[Answer to Q8.5]
\textbf{The obligation is real, forward-looking, and distinct from incident response --- and it has a named owner here.} The supply-chain-security duty of \loc{cybersecurity-baseline-regime} (in the EU, \gls{nis2} Article~21(2)(d)~\cite{nis2}) requires ongoing management of supplier and supply-chain security, not merely post-incident remediation.

\medskip\noindent\textbf{The architecture already carries the control.} The lifecycle wave, a supply-chain-attestation instantiation (signed provenance and SBOMs), and a procurement-clause workstream established in the foundations wave~\cite{ng-tech-ref-2026} produce software provenance and contractual security clauses applied to \emph{every} acquisition, regardless of which zone procures it.

\medskip\noindent\textbf{One accountable owner, mirroring the Article~23 clock.} Just as the central \gls{soc} owns the \gls{nis2} reporting clock (Q8.1), a central security-governance owner maintains the supplier-assurance baseline --- the approved-supplier register, the mandatory procurement clauses, and the attestation requirements --- while each zone charter binds local procurement to that baseline. Federation therefore does not fragment supply-chain assurance: it gives it a single, locatable owner with per-zone obligations reporting inward.
\end{answerbox}

\bigskip

% --- Q8.6 ---
\begin{questionbox}[Q8.6 --- ``Data-subject rights are unfulfillable across federated zones and a distributed audit store'']
Q8.3 and Q8.4 establish the lawful basis, DPIA and proportionality of monitoring, but never explain how the institution actually fulfils a data subject's Article~15 access, Article~16 rectification or Article~17 erasure request when identity is federated and the audit log is distributed across autonomous zones. A \gls{dpo} could refuse on the ground that rights become unfulfillable in a distributed design.
\end{questionbox}

\begin{answerbox}[Answer to Q8.6]
\textbf{The concern is well founded, and the same single-accountable-owner logic that locates the Article~23 clock (Q8.1) locates rights fulfilment.} Distribution raises a genuine operability question; the design answers it with one front door, one map, and per-zone duties.

\medskip\noindent\textbf{Central intake (one front door).} A data-subject request is received and tracked at a single \gls{dpo}-owned intake point, so the data subject never has to address individual zones.

\medskip\noindent\textbf{Per-zone retrieval duty.} Each zone charter carries a binding obligation to locate, return, rectify or erase the records it holds within a sub-deadline, reporting inward to the \gls{dpo} --- exactly as zones report inward to the \gls{soc} for incident notification.

\medskip\noindent\textbf{Bounded retention and a locator.} The audit pipeline records minimised \emph{metadata} for a bounded period (Q8.4), so scheduled expiry makes Article~17 erasure tractable rather than a distributed deletion hunt; and the Article~30 Record of Processing Activities~\cite{gdpr}, owned by the \gls{dpo}, maps which zones and systems hold which categories of data, so any Article~15/16/17 request is routed deterministically across the federation.

\medskip\noindent\textbf{The evidence store is itself a control.} The audit trail that underwrites these rights is not assumed trustworthy: it is append-only and tamper-evident (per-record signing or hash-chaining), with separation of duties between operators and readers and a four-eyes rule on retention expiry, as specified in the Technical Reference's observability chapter. The record proving accountability therefore cannot be silently altered or selectively deleted~(\gls{gdpr} Art.~32 / Art.~5(1)(f); ISO/IEC~27001 A.8.15).
\end{answerbox}

% ============================================================================
\section{Questions on costs and feasibility}
\label{sec:q-costs}
% ============================================================================

\epigraph{``There is nothing more expensive than a missed opportunity.''}{--- Attributed to H.\ Jackson Brown Jr.}

% --- Q9.1 ---
\begin{questionbox}[Q9.1 --- ``The cost analysis lacks concrete numbers --- how much will this actually cost?'']
The dossier discusses costs qualitatively but does not provide a detailed budget. Without concrete figures, the cost-benefit analysis is speculative.
\end{questionbox}

\begin{answerbox}[Answer to Q9.1]
\textbf{This is a fair request} and one that can only be fully addressed through a detailed implementation study specific to the institution's current infrastructure, contracts, and staffing.

\medskip\noindent\textbf{What the dossier does provide:}
\begin{itemize}[nosep]
    \item \textbf{The governance step} (commission, charters, digital-policy commission): predominantly \emph{organisational} --- no new infrastructure, implementable with existing staff;
    \item \textbf{The pilot step} (selected entities with domain expertise): standard open-source tools and \gls{nren} services already available through \gls{geant}~\cite{geant};
    \item \textbf{The extension step:} gradual rollout aligned with validated results;
    \item \textbf{Cost of inaction:} the global all-industry average data-breach cost reached \$4.88M in 2024~\cite{ibm-cost-breach-2024} (the education sector sits below this average); talent drain, Shadow IT, and lost competitiveness in funding calls are additional indirect costs.
\end{itemize}

\noindent The absence of a line-item budget in a \emph{policy proposal} is normal --- the budget is produced during the implementation study that follows the governance step. The purpose is to demonstrate that the approach is \emph{architecturally sound and scalable from a teaching-led to a research-intensive institution}, not to substitute for a procurement exercise.

\smallskip\noindent\textit{\footnotesize The standalone Technical Reference~\cite{ng-tech-ref-2026} provides indicative operational envelopes per layer, three reference pace options for the migration (5/7/10-year horizons), and explicit discussion of the financial-sovereignty dimension of every architectural choice.}
\end{answerbox}

\bigskip

% --- Q9.2 ---
\begin{questionbox}[Q9.2 --- ``The transition plan is too optimistic'']
Institutional change at universities is notoriously slow. The proposed progressive transition may face delays due to institutional inertia.
\end{questionbox}

\begin{answerbox}[Answer to Q9.2]
\textbf{The progressive structure is designed to absorb delays:}
\begin{itemize}[nosep]
    \item NC~State completed a governance redesign within a comparable scope~\cite{ncsu-governance};
    \item NSF CC* grants that fund \gls{sciencedmz} deployments typically run two to three years for full deployment at a single campus~\cite{nsf-cc};
    \item The governance step is primarily \emph{commission, charters, and documents} --- it does not require infrastructure procurement, the usual bottleneck.
\end{itemize}

\noindent\textbf{The transition is incremental, not deadline-driven.} The progressive structure ensures that delays in one step do not create systemic risk --- the institution benefits incrementally from each completed step.

\medskip\noindent\textbf{The alternative timeline is worse:} the current state (monolithic restriction, Shadow IT, regulatory uncertainty) does not improve with time. Every month of delay is a month of talent drain, Shadow IT risk, and missed compliance opportunities.
\end{answerbox}

\bigskip

% --- Q9.3 ---
\begin{questionbox}[Q9.3 --- ``The economic case is asserted, not shown --- the board cannot decide without figures'']
This dossier asks the governing bodies to adopt the policy; the detailed, situation-specific costing would properly be commissioned from the CIO once they decide to proceed. But a board still needs an order of magnitude to take that decision --- and as written the dossier gives no figures, so a CIO could note that the economic case has not been framed.
\end{questionbox}

\begin{answerbox}[Answer to Q9.3]
\textbf{The detailed, situation-specific cost analysis is the CIO's to produce once the governing bodies request it; the dossier's task is to give the board a defensible order of magnitude and a cost shape, not to pre-empt that analysis.}
\begin{enumerate}[nosep]
  \item \textbf{Cost shape.} The governance step (commission, charters, digital-policy commission) is \emph{cost-neutral} --- existing staff, no procurement. Cost-bearing items sit in later, optional waves (segmentation, Science~DMZ, observability storage), each gated by the validated results of the prior step.
  \item \textbf{A scalable benchmark.} Rather than invent a figure, the order of magnitude can be read from sector benchmarks: higher education currently allocates only about 3.6\% of the IT budget to information security~\cite{educause-spending}, below the cross-industry averages reported in vendor benchmark data~\cite{gartner-security-spend2024}. A board can therefore size the likely envelope from its enrolment and current IT budget, and ask the CIO to refine it.
  \item \textbf{Cost of inaction.} The status quo is not free: the global all-industry average data-breach cost reached \$4.88M in 2024~\cite{ibm-cost-breach-2024} (the education sector below it), with talent drain and Shadow-IT exposure growing with delay.
\end{enumerate}
\noindent This frames precisely the decision the dossier asks of the governing bodies --- to commission a costed implementation study --- with enough quantified context (a scalable percentage and a quantified downside) to make that decision defensible, while leaving the precise TCO to the CIO.
\end{answerbox}

\bigskip

% --- Q9.4 ---
\begin{questionbox}[Q9.4 --- ``Even if it is affordable, a small IT team cannot operate a heavyweight open-source stack'']
The cost objection has a sharper, operational form: the reference stack names heavyweight platforms whose day-to-day operation demands specialist staff a small or teaching-led institution does not have. Affordability of acquisition is not the same as feasibility of operation.
\end{questionbox}

\begin{answerbox}[Answer to Q9.4]
\textbf{The reference stack is the research-intensive instantiation, not a mandatory floor; the architecture is defined by its functional layers and the sovereignty grid, and each layer degrades to a managed or mutualised realisation as scale decreases.}
\begin{enumerate}[nosep]
  \item \textbf{Operate by responsibility, not by platform.} A down-scaled profile draws compute and storage from a national research cloud, an \gls{nren} service, or a managed provider under a data-processing agreement~\cite{ng-tech-ref-2026}: the institution retains sovereignty over \emph{where} and \emph{under what terms} data sits without operating the underlying platforms itself.
  \item \textbf{Roles are not headcount.} The role inventory describes responsibilities that, in a small institution, several of which are discharged by the same person; it is not summed into a minimum specialist team.
  \item \textbf{Optional layers stay optional.} Components justified only at scale --- the Science~DMZ, dedicated observability storage --- can be deferred entirely without weakening the rest of the architecture.
\end{enumerate}
\noindent The operational burden therefore scales with the chosen instantiation, not with the architecture: a small institution adopts the same principles at a managed-service realisation it can actually run.
\end{answerbox}

\bigskip

% --- Q9.5 ---
\begin{questionbox}[Q9.5 --- ``Our cyber-insurer's questionnaire forbids this: EDR everywhere, no unmanaged admin rights, full perimeter inspection --- adopting the proposal voids our coverage'']
Typical underwriting questionnaires require endpoint detection on
every host, prohibit unmanaged administrator rights, and assume
perimeter inspection of all traffic. Accountable researcher root
access and a \gls{sciencedmz} that bypasses the enterprise firewall
facially contradict all three. Whatever the architecture's merits,
the institution cannot operate uninsured.
\end{questionbox}

\begin{answerbox}[Response Q9.5]
\textbf{This concern has merit in its premise:} insurability is a
genuine external gate, and ``the architecture is sound'' is not an
answer an underwriter accepts. The response, however, runs in three
steps --- and its order matters.

\begin{enumerate}[nosep]
  \item \textbf{The proposal strengthens the insurable posture.}
  Network segmentation, institution-wide \gls{mfa}, federated
  identity and classification-driven observability are precisely
  the control families underwriting commonly rewards; researcher
  root access is \emph{accountable}, not unmanaged --- charter-bound,
  logged, revocable, with a named risk owner per zone (the enclave
  charter); and the \gls{sciencedmz} is disclosable as a
  compensating-control architecture (ACL+IDS profile, dedicated
  transfer nodes, no general-purpose hosts --- the documented ESnet
  pattern). The per-zone evidence of P14 and the maturity reporting
  produce, as a by-product, the documentation underwriters ask for
  at renewal.

  \item \textbf{Where the questionnaire still conflicts, the first
  reflex is to renegotiate the contract, not the mission.} The
  report's \keyword{liberty-first rule} is that constraint is never
  preferred to freedom by default: a contract that would suppress an
  academic freedom is a contract to renegotiate. Concretely: engage the
  insurer or broker at Wave~0 with the zone catalogue and charters
  in hand; declare deviations as documented exceptions with
  compensating controls rather than letting a generic corporate
  questionnaire default the institution to ``non-compliant''; and
  seek academic-context underwriting terms. Freedom has a price the
  governance body can knowingly choose to pay --- retention of a
  residual risk, an alternative carrier, or sector mutualisation ---
  rather than letting a standard form rewrite the institution's
  mission.

  \item \textbf{At sector scale, the response is collective.} Where
  imposed terms would make academic digital freedom uninsurable by
  default, universities act together --- through their consortia,
  NREN communities and sector bodies, in the same spirit as the
  dossier's Position~C (influence and reform): engaging insurers,
  and where necessary the bodies that regulate them, so that the
  insurability of a university is assessed against the controls it
  actually operates, not against a template written for a
  single-perimeter corporation.
\end{enumerate}

\noindent\textbf{Bottom line.} The questionnaire is an input to a
negotiation, not a design authority. The proposal gives the
institution both better terms to ask for and the evidence to ask
with; the truly unanswerable position is the undocumented flat
network.
\end{answerbox}

% ============================================================================
\section{Questions on methodology, process and credibility}
\label{sec:q-meta}
% ============================================================================

\epigraph{``The important thing is not to stop questioning. Curiosity has its own reason for existing.''}{--- Albert Einstein, \textit{LIFE Magazine}, 1955}

% --- Q10.1 ---
\begin{clarificationbox}[Q10.1 --- ``This document is biased --- it was written to support a predetermined conclusion'']
The dossier reads as an advocacy document, not an objective analysis. It selects evidence that supports federation and ignores evidence that supports centralisation.
\end{clarificationbox}

\begin{answerbox}[Answer to Q10.1]
\textbf{Every policy document has a thesis.} The distinction between a biased document and a well-argued one is whether the argument is supported by evidence and whether counter-arguments are addressed.

\medskip\noindent\textbf{The evidence base:}
\begin{itemize}[nosep]
    \item An extensive, fully referenced bibliography with verified sources;
    \item multiple reference university case studies across several countries;
    \item Major normative frameworks (NIST, ISO, \gls{cobit}, EU regulations);
    \item Sector organisations (EDUCAUSE, \gls{geant});
    \item Independent sector reports (Verizon DBIR, Sophos);
    \item A body of court decisions and legal instruments.
\end{itemize}

\noindent\textbf{Counter-arguments are addressed:} the incident analysis (Q7.1, Q7.2) directly engages with the strongest evidence that could be used \emph{against} the thesis and demonstrates that it actually supports it. The risk register honestly catalogues the risks of the proposed approach. This very module demonstrates willingness to engage with criticism.

\medskip\noindent It is worth noting that maintaining the status quo is also a policy position that benefits from the same evidence-based scrutiny. The status quo is itself a position that should be defended with evidence --- not treated as the neutral default.
\end{answerbox}

\bigskip

% --- Q10.2 ---
\begin{clarificationbox}[Q10.2 --- ``This is an attack on the IT department / on the security function personally'']
The dossier criticises the current approach so extensively that it amounts to a vote of no confidence in the IT department's competence.
\end{clarificationbox}

\begin{answerbox}[Answer to Q10.2]
\textbf{The dossier critiques an \emph{approach}, not \emph{people}.} The monolithic restriction approach is common in many institutions and is often adopted in good faith by competent professionals responding to genuine threats. The argument is that a \emph{different} approach is more effective for a \emph{university context} --- not that the IT department is incompetent.

\medskip\noindent\textbf{The proposals benefit the IT department:}
\begin{itemize}[nosep]
    \item \emph{Reduced perimeter:} the IT department focuses on the administrative enclave rather than trying to manage everything;
    \item \emph{End of ad hoc exceptions:} charters and the \gls{icradp} replace informal, time-consuming negotiations;
    \item \emph{Clear accountability:} departments take formal responsibility for their zones instead of blaming central IT when things go wrong;
    \item \emph{Better regulatory posture:} an evidence-producing architecture strengthens the security function's position in audits.
\end{itemize}

\noindent Policy evolution is a normal institutional process. Every institution periodically re-evaluates its approaches; this dossier is a contribution to that ongoing dialogue, not an indictment.
\end{answerbox}

\bigskip

% --- Q10.3 ---
\begin{clarificationbox}[Q10.3 --- ``The document was partially written with AI tools --- therefore it cannot be trusted'']
The dossier acknowledges using AI tools in its composition. This undermines its credibility as a serious academic contribution.
\end{clarificationbox}

\begin{answerbox}[Answer to Q10.3]
\textbf{This is a \emph{genetic fallacy}:} evaluating an argument by its origin rather than its content. The validity of the dossier depends on:
\begin{enumerate}[nosep]
    \item The accuracy of its factual claims (verifiable through its extensive base of cited sources);
    \item The soundness of its reasoning (evaluable on its logical merits);
    \item The robustness of its recommendations (assessable against established frameworks).
\end{enumerate}

\noindent None of these are affected by whether the text was drafted with AI assistance, a typewriter, or a quill pen. The \gls{aiact} itself recognises this: it regulates AI \emph{systems deployed in the real world}, not AI used as a writing tool~\cite{eu-ai-act}. LLMs produce nothing without human direction; the differentiating factor is the disciplined, review-bound workflow around the tool, not the tool itself.

\medskip\noindent\textbf{Transparency:} the explicit disclosure of AI-assisted composition is an act of intellectual honesty, not a weakness. Demanding disclosure and then using the disclosure to discredit the work is a rhetorical trap.

\medskip\noindent\textbf{Scientific standard:} the relevant question for any document is whether the claims are accurate, the sources verifiable, and the reasoning sound. These criteria apply regardless of the drafting tools used.
\end{answerbox}

% ============================================================================
\section{Summary: quick-reference table}
\label{sec:summary}
% ============================================================================

\begin{longtable}{>{\raggedright}p{1cm} >{\raggedright}p{5.5cm} >{\raggedright\arraybackslash}p{7.5cm}}
\caption{Quick-reference summary of questions and key points} \label{tab:arg-summary} \\
\toprule
\textbf{Ref.} & \textbf{Question (summary)} & \textbf{Key points} \\
\midrule
\endfirsthead
\toprule
\textbf{Ref.} & \textbf{Question (summary)} & \textbf{Key points} \\
\midrule
\endhead
\midrule
\multicolumn{3}{r}{\footnotesize\itshape Continued on next page} \\
\endfoot
\bottomrule
\endlastfoot

Q1.1 & Segmentation increases complexity & Complexity $\neq$ attack surface; NIST mandates segmentation; phased rollout + SDN mitigate \\
\addlinespace
Q1.2 & Single perimeter = better visibility & Single firewall creates bottlenecks; per-zone monitoring gives greater total observability \\
\addlinespace
Q1.3 & ``You just want less security'' & \emph{Ad hominem}; proposals endorsed by NIST, ISO, EDUCAUSE, and a broad set of reference universities \\
\addlinespace
Q1.4 & Firebreak analogy is misleading & Analogy has limits; principle stands on NIST/Verizon evidence independently \\
\addlinespace
Q1.5 & Zero Trust claimed but SASE/ZTNA excluded & Maps to all seven SP~800-207 tenets; SASE/ZTNA declined as a sovereignty choice, not Zero Trust \\
\addlinespace
Q2.1 & Federation creates accountability gaps & Charters + SOC + zone owners create \emph{more} granular accountability \\
\addlinespace
Q2.2 & Departments will be irresponsible & Risk acknowledged; mitigated by charters, SOC, default-deny, isolation, audits \\
\addlinespace
Q2.3 & Undermines the security function's authority & COBIT/Weill-Ross: governance $\neq$ management; no major university gives one officer unilateral authority \\
\addlinespace
Q2.4 & Academics lack security expertise & Distinction between technical implementation and policy direction; every framework requires stakeholder input \\
\addlinespace
Q3.1 & Root access is unacceptable risk & ETH model: accountable root access + isolation is safer than shadow root access \\
\addlinespace
Q3.2 & Autonomy unfairly burdens researchers & Central IT still provides backbone/SOC; Data Stewards + local ISGs bridge the gap \\
\addlinespace
Q3.3 & Science DMZ bypasses security & Replaces firewall with \gls{acl}s + Zeek \gls{ids}; endorsed by DoE/NSF; 100+ deployments \\
\addlinespace
Q3.4 & Not deployed locally, so unsuitable & Appeal to absence; substrate often present; non-deployment reflects funding routes, not merit; mutualised form proposed \\
\addlinespace
Q4.1 & BYOD is inherently insecure & BYOD exists regardless; progressive interoperability + Zero Trust is safer than unmanaged \\
\addlinespace
Q4.2 & Commercial AI creates \gls{gdpr} liability & Differentiated framework: adequacy or \gls{scc}s+\gls{tia} with supplementary measures, data classification, AI Act risk-based classification, traceability \\
\addlinespace
Q4.3 & AI access is convenience, not necessity & Status-quo bias; AI is structural for research/teaching; blocking creates shadow AI \\
\addlinespace
Q5.1 & We are not MIT/Stanford & False equivalence; principles are scale-independent; evidence spans institutions of all sizes \\
\addlinespace
Q5.2 & Cherry-picked success stories & Broad set of reference universities + sector consensus (EDUCAUSE); counter-examples welcome \\
\addlinespace
Q5.3 & Cannot retrofit federation & NC State, Penn State, Wisconsin-Madison, Michigan deliberately transitioned to federation \\
\addlinespace
Q5.4 & Reversibility rests on single-vendor tests & Exit tests are cross-implementation by design; tested behaviour, not vendor promise \\
\addlinespace
Q6.1 & Adjustments commission creates bureaucracy & Short published deadlines; handles exceptions, not routine; formalises existing informal process \\
\addlinespace
Q6.2 & Commission could be misused & Structural safeguards: documented justification, proportionality, security function heard, transparency \\
\addlinespace
Q6.3 & Citing rights is disproportionate & Binding EU/ECHR law, not philosophy; CJEU/ECtHR case law directly applies \\
\addlinespace
Q6.4 & Technostress is pseudoscience & 105-study systematic review (La Torre 2019); peer-reviewed in top journals \\
\addlinespace
Q7.1 & Recent breaches prove need for restriction & Root causes are supply chain, credentials, lack of segmentation --- not openness \\
\addlinespace
Q7.2 & Incident analysis is biased & Interpretation aligns with Verizon DBIR, NIST, Sophos independently \\
\addlinespace
Q8.1 & NIS2 requires centralisation & NIS2 is requirements-based, not architecture-prescriptive; federation meets all requirements with one clock owner \\
\addlinespace
Q8.2 & Risk-based argument is too convenient & Applies symmetrically; state of the art is segmented architecture; EDPB/ENISA confirm \\
Q8.3 & No \gls{gdpr} lawful basis / DPIA trigger stated & Art.~6(1)(e)/(f) basis stated; Art.~35 DPIA expected; segmentation aids data minimisation \\
Q8.4 & B\u{a}rbulescu not applied to own monitoring & Metadata-only, per-zone, disclosed monitoring satisfies the six criteria \\
Q8.5 & NIS2 supply-chain (Art.~21(2)(d)) unowned & Attestation + procurement clauses; central assurance owner, per-zone duties \\
Q8.6 & Data-subject rights unfulfillable when federated & Central \gls{dpo} intake + per-zone retrieval duty + bounded retention + Art.~30 RoPA locator \\
\addlinespace
Q9.1 & No concrete budget numbers & Normal for policy documents; budget produced after governance step; open-source/NREN tools minimise cost \\
\addlinespace
Q9.2 & Timeline is too optimistic & Phased structure absorbs delays; status quo worsens with time \\
\addlinespace
Q9.3 & Economic case is asserted, not shown & Cost shape + scalable benchmark (higher education's 3.6\% vs cross-industry averages) + cost of inaction frame the board decision \\
\addlinespace
Q9.4 & Small team cannot operate the stack & Reference stack is the high-scale instantiation; layers degrade to managed/mutualised realisations \\
\addlinespace
Q9.5 & Insurer's questionnaire forbids researcher root / Science DMZ & Underwriting credits + disclosable compensating controls; renegotiate for academic-context terms (questionnaire is negotiation input, not design authority); collective action in the spirit of Position~C \\
\addlinespace
Q10.1 & Document is biased & Extensively referenced, counter-arguments addressed, risk register honest \\
\addlinespace
Q10.2 & Attack on the IT department & Critiques approach, not people; proposals benefit IT department \\
\addlinespace
Q10.3 & AI-assisted = untrustworthy & Genetic fallacy; validity depends on accuracy/reasoning, not drafting tools \\

\end{longtable}

% ============================================================================
\section{Summary and outlook}
\label{sec:arg-conclusion}
% ============================================================================

\epigraph{``The test of a first-rate intelligence is the ability to hold two opposed ideas in mind at the same time and still retain the ability to function.''}{--- F.\ Scott Fitzgerald, \textit{The Crack-Up}, 1936}

The systematic treatment of the frequently raised questions above supports four conclusions:

\begin{enumerate}
    \item \textbf{No question invalidates the core thesis.} Every genuine concern (complexity, accountability, costs, feasibility) is met by a specific mechanism (phased rollout, charters, \gls{soc} oversight, open-source tooling). Some questions have partial merit and highlight implementation risks that must be managed; the risk register acknowledges these.

    \item \textbf{Several questions rest on premises that, on examination, require correction or additional context.} These are addressed through clarification rather than rebuttal, and the underlying concerns are taken seriously in each case.

    \item \textbf{The evidence base is robust:} an extensive, fully referenced base, multiple documented university case studies, and alignment with NIST, ISO, \gls{cobit}, the EU regulatory framework, EDUCAUSE, and independent sector reports. Engaging constructively with this evidence is the most productive path toward sound IT governance decisions.

    \item \textbf{The strongest argument is convergence:} MIT, Stanford, Cambridge, Oxford, ETH~Z\"urich, CERN, Michigan, Wisconsin-Madison, NC~State, Penn State --- and the 100-plus universities operating a \gls{sciencedmz} --- converge on segmented, federated architectures, while NIST, ISO, EDUCAUSE, \gls{geant}, and the EU regulatory framework all support risk-based, proportionate security. When both practice and theory converge this strongly, an institution choosing a fundamentally different path should articulate the evidence supporting that alternative.
\end{enumerate}

\medskip\noindent The proposals are grounded in extensive evidence, aligned with major governance frameworks, and internationally validated. Their implementation will require careful management, but they offer a robust foundation for constructive dialogue. Further questions are welcome.
